\documentclass[11pt]{article}

\usepackage[margin=1.1in]{geometry}
\usepackage{booktabs}
\usepackage{amsmath}
\usepackage[hidelinks]{hyperref}
\usepackage{parskip}

\newcommand{\code}[1]{\texttt{#1}}

\sloppy
\emergencystretch=3em

\title{Gang placement under a greedy allocator:\\
feasibility cliff and the anchor heuristic}
\author{medios research track}
\date{June 9, 2026}

\begin{document}

\maketitle

\begin{center}
\small Companion to \code{docs/research/offers.md}; experiment tracked as
medios issue \#316 (private history).
Simulation source: \code{tests/sim\_gang\_placement.cpp}.
\end{center}

\begin{abstract}
A gang reservation is feasible only if all members place jointly, so
feasibility is a conjunction over members and the architecture question for
medios is whether a greedy allocator can carry gangs, or whether they force
global coordination. We test this in a spatial-only placement simulation: no
time dimension, no medios code, synthetic clusters of labeled nodes under
packed background load. Three placement strategies run on identical
scenarios: in-order first-fit (\emph{naive}), most-constrained-first with
backtracking only on the most constrained member (\emph{anchor}), and a
best-fit-decreasing reference (\emph{bfd}) validated as exactly deciding
feasibility on this demand mix by exhaustive search over 5{,}000 small
instances. Results: the multiplicative feasibility cliff is real --- naive
feasibility falls from 1.0 to 0.49 between gang sizes 8 and 64 at 50\%
utilization with capacity to spare. The anchor strategy matches the exact
reference in 71 of 72 sweep cells to four decimals (one divergent trial in
144{,}000), so greedy placement with the right ordering does not
approximate the feasibility frontier --- it sits on it. Constraint
correlation between members steepens the cliff and widens the gap over
naive. One claim is refuted as stated: anchoring does not cut placement
effort under honest accounting; its feasibility-set precomputation costs
about six times naive's node examinations, while bfd achieves the same
feasibility with neither precomputation nor backtracking. The decision the
experiment was built to make: medios can extend the greedy bin-packer for
gangs; ordering, not global search, is the missing ingredient.
\end{abstract}

\section{Hypothesis}

From \code{offers.md}: gang feasibility is a conjunction over members
dominated by the member with the smallest feasible set. Three testable
claims:

\begin{itemize}
\item \textbf{H1 (cliff).} Joint feasibility decays multiplicatively in
  gang size $N$, even when aggregate capacity is plentiful. The member
  count itself is a feasibility cliff.
\item \textbf{H2 (anchor).} Placing the smallest-feasible-set member first
  and backtracking only on it keeps a greedy allocator near the optimal
  feasibility frontier and reduces placement effort.
\item \textbf{H3 (predictor).} The size of the smallest feasible set at
  scenario start predicts the outcome.
\end{itemize}

The refutation criteria fixed in advance: anchor $\approx$ naive would mean
gangs need global coordination; no cliff would mean the fragility argument
is overstated.

\section{Experiment}

\textbf{Model.} $M = 100$ nodes, 8 slots each, one binary topology label
(gpu type, 50/50). A gang member demands $d \in \{1,2,4,8\}$ slots
(uniform) on a single node with matching label. Background load is itself a
packed workload: tasks with the same size mix placed first-fit in random
node order until target utilization $U$ is reached. Gang members optionally
correlate: with probability $\rho$ a member requires the gang's home type
instead of an independent draw, parameterizing the independence assumption
behind the $0.9^N$ argument ($\rho = 0$ independent, $\rho = 1$ whole gang
in contention for one node population).

\textbf{Strategies}, paired on identical scenarios with a fixed seed:

\begin{itemize}
\item \emph{naive} --- members in arrival order, first-fit over nodes in
  index order, fail on the first stuck member.
\item \emph{anchor} --- members ordered by feasible-set size ascending; the
  most constrained member is the only backtrack point and is retried across
  its feasible nodes; the rest place greedily.
\item \emph{bfd} --- best-fit-decreasing by demand, no backtracking; the
  reference. The demand sizes are divisible ($1 \mid 2 \mid 4 \mid 8$), for
  which best-fit-decreasing is expected to decide feasibility exactly;
  rather than assume the theorem, \code{--selfcheck} compares bfd against
  exhaustive search on 5{,}000 random small instances: zero mismatches.
\end{itemize}

\textbf{Methodological note.} The first background-load model (per-slot
binomial occupancy) degenerated the experiment: it leaves almost no fully
free nodes at any utilization, so every whole-node member was infeasible at
scenario start and no ordering could matter. The packed-workload model was
adopted because failures under it are caused by competition and
fragmentation --- the thing under test --- rather than by an artifact of
the load generator.

\textbf{Sweep.} $N \in \{2,4,8,16,24,32,48,64\}$ $\times$
$U \in \{0.50, 0.70, 0.85\}$ $\times$ $\rho \in \{0, 0.5, 1\}$, 2{,}000
trials per cell. Reproduce with:
\code{./sim\_gang\_placement --trials 2000 --seed 42}.

\section{Results}

\begin{table}[h]
\centering
\begin{tabular}{lccc}
\toprule
$N$ & naive & anchor & bfd (exact) \\
\midrule
8  & 1.000 & 1.000 & 1.000 \\
16 & 0.998 & 0.998 & 0.998 \\
32 & 0.954 & 0.968 & 0.968 \\
48 & 0.788 & 0.830 & 0.830 \\
64 & 0.488 & 0.596 & 0.596 \\
\bottomrule
\end{tabular}
\caption{Feasibility rate at $U = 0.50$, $\rho = 0$ (independent members).}
\end{table}

\begin{table}[h]
\centering
\begin{tabular}{lccc}
\toprule
$\rho$ & naive & anchor & bfd (exact) \\
\midrule
0   & 0.788 & 0.830 & 0.830 \\
0.5 & 0.626 & 0.712 & 0.712 \\
1   & 0.251 & 0.375 & 0.375 \\
\bottomrule
\end{tabular}
\caption{Effect of member correlation at $N = 48$, $U = 0.50$.}
\end{table}

\begin{itemize}
\item \textbf{H1 confirmed.} The cliff is multiplicative and appears while
  aggregate capacity is ample (at $N = 64$, $U = 0.50$ the gang demands
  ${\sim}240$ of ${\sim}400$ free slots and naive still fails half the
  time). Higher utilization pulls the cliff to smaller $N$; at $U = 0.85$
  feasibility is starvation-dominated (members with empty feasible sets),
  where no ordering helps.
\item \textbf{H2 confirmed on feasibility, refuted on effort.} Anchor
  equals the exact reference in every cell except one divergent trial in
  144{,}000 ($N = 64$, $\rho = 1$: 0.0595 vs 0.0600). Greedy with
  most-constrained-first ordering is not approximately good --- it is
  indistinguishable from optimal in this model. On effort the claim
  inverts: anchor examines ${\sim}6\times$ more nodes than fail-fast naive
  because the feasible-set ranking costs $N \times M$ examinations up
  front. bfd reaches the same feasibility with no precomputation and zero
  backtracks.
\item \textbf{H3 partially confirmed.} Mean smallest-feasible-set size
  separates successes from failures in every cell (e.g.\ 9.5 vs 7.6 at
  $N = 64$, $U = 0.50$), but the separation narrows as $N$ grows; it is a
  useful signal, not a sharp threshold.
\item \textbf{Correlation steepens the cliff and widens the anchor
  advantage} (at $N = 48$, $\rho = 1$ the relative gain over naive is
  ${\sim}50\%$). The independence assumption in the $0.9^N$ argument is
  optimistic; real gangs that compete for one resource population are more
  fragile than the independent model predicts, and ordering matters more
  for them, not less.
\end{itemize}

\textbf{Why anchor collapses onto bfd here.} With a single label axis,
feasible-set size is nearly monotone in demand, so ``smallest feasible set
first'' and ``largest demand first'' induce the same order. The model
cannot yet distinguish the anchor heuristic from plain demand ordering.
They separate only under multi-axis constraints (e.g.\ rack $\times$ type),
where a small demand can still have a small feasible set.

\textbf{Threats to validity.} One label axis and two values; single-node
members only (no member spanning); demand sizes divisible, which both makes
bfd exact and flatters all packers; $M = 100$ fixed; no time dimension by
design. The result decides the v0 question only.

\section{Conclusion}

The decision criterion is met. Gang placement does not need global
coordination in the spatial regime: a greedy allocator with
most-constrained-first ordering sits on the exact feasibility frontier, and
even the cheap demand-descending variant does. The cliff is real,
correlation makes it worse, and the practical content of
``anchor-on-the-bottleneck'' at this stage is an ordering rule, not a
search procedure. This supports extending the existing medios bin-packer
for gang reservations and keeping placement policy negotiable by frameworks
(FAC3-0014), rather than adopting a solver.

\section{Next steps}

\begin{enumerate}
\item \textbf{Separate anchor from demand ordering (v0.5).} Add a second
  constraint axis (rack $\times$ type) so feasible-set size and demand
  decouple; measure whether anchor's advantage over bfd appears where the
  theory says it must.
\item \textbf{Interval dimension (v1).} Node calendars and member windows:
  test the placement-cost vs calendar-cost split and the ``longest window
  fragments compaction'' claim from \code{offers.md}.
\item \textbf{Non-divisible demands.} Break the divisibility property and
  measure how far bfd and anchor fall from exhaustive feasibility; bounds
  the optimism of this model.
\item \textbf{medios integration (v2).} Only after v1: gang offers through
  the real allocator, composing with the reservation primitive verified in
  \code{reservation-mvp.md}.
\end{enumerate}

\end{document}
